% templates/bifurcation_spectrum.tex.mustache
% AUTO-GENERATED Stage 5 bifurcation panel
% Rendered by scripts/build_campaign_report.jl
% DO NOT EDIT MANUALLY

\section{Physical Interpretation of the Continuation Parameter}

The continuation parameter $\gamma$ represents the effective atmosphere--surface coupling strength inferred from the sparse structural operator identified in Stage~4.
Small values of $\gamma$ correspond to weak turbulent transport and strongly decoupled nocturnal conditions, where the surface inversion evolves in relative isolation from the overlying residual layer.
Larger values of $\gamma$ represent enhanced vertical exchange between the surface inversion and the residual layer above, sustaining coherent turbulent structures through the column.

Parameter continuation therefore traces the geometric deformation of the reconstructed attractor manifold as the atmospheric state transitions between coupled and decoupled regimes.
The leading eigenvalue $\alpha(\gamma) = \max\,\Re(\lambda_i)$ serves as a scalar indicator of manifold stability: negative values indicate a contracting, stable attractor; the zero crossing at $\gamma_c$ marks the onset of structural instability; and the imaginary part $\beta(\gamma) = |\mathrm{Im}(\lambda_\mathrm{dom})|$ tracks the intrinsic oscillatory frequency of the emerging mode.

\section{Stage 5 Bifurcation Spectrum}

This panel summarizes the continuation branch spectral abscissa $\alpha(\gamma)$ for campaign \texttt{ BLLAST }.


\begin{figure}[htbp]
    \centering
    \usepgfplotslibrary{groupplots}
    	tikzsetnextfilename{bifurcation-panel-bllast}
    \begin{tikzpicture}
        \begin{groupplot}[
            group style={
                group name=bifurcation_plots,
                group size=1 by 2,
                vertical sep=12pt,
                x descriptions at=edge bottom
            },
            width=0.96\textwidth,
            grid=both,
            grid style={line width=.1pt, draw=gray!15},
            major grid style={line width=.2pt, draw=gray!35},
            xmin=0.001000, xmax=1.000000,
            tick label style={font=\small},
            label style={font=\small},
            axis background/.style={fill=gray!2}
        ]

        \nextgroupplot[
            height=4.5cm,
            ylabel={Growth $\alpha = \max\,\Re(\lambda)$},
            axis lines=middle,
            axis line style={draw=gray!60},
            legend style={at={(0.98,0.05)}, anchor=south east, font=\small},
            xticklabels={}
        ]
            \addplot[
                color=blue!70!black,
                line width=1.3pt,
                mark=none
            ] table[
                col sep=comma,
                x=gamma,
                y=max_real_eig
            ] {../../data/outputs/stage5_bifurcation_branches_bllast.csv};
            \addlegendentry{$\alpha(\gamma)$}


        \nextgroupplot[
            height=3.8cm,
            xlabel={Dissipation parameter ($\gamma$)},
            ylabel={Frequency $\beta\,(\mathrm{rad\,s^{-1}})$},
            ymin=0.0, ymax=0.029704,
            axis lines=left,
            axis line style={draw=gray!60},
            legend style={at={(0.98,0.92)}, anchor=north east, font=\small}
        ]
            \addplot[
                color=teal!80!black,
                line width=1.3pt,
                mark=none
            ] table[
                col sep=comma,
                x=gamma,
                y=max_imag_eig
            ] {../../data/outputs/stage5_bifurcation_branches_bllast.csv};
            \addlegendentry{$\beta(\gamma)$ frequency branch}

        \end{groupplot}
    \end{tikzpicture}
    \caption{Stage 5 parameter continuation spectrum: the top panel charts the principal manifold stability branch $\alpha(\gamma)$, while the bottom panel shows the aligned oscillation-frequency branch $\beta(\gamma)$ when available.}
    \label{fig:bifurcation_spectrum_bllast}
\end{figure}




\noindent\textit{Cross-campaign Stage 5 comparison table omitted because one or more required summary manifests were not available at render time.}
